\section{Critical boundedness and the nonuniqueness interval}
\label{con:section}

We combine the preceding estimates to prove the critical operator bound.
The argument fixes all path constants before choosing the subcritical
parameter and the deformation time. We then transfer the bound to an
interval above criticality and prove the simultaneous cluster assertion.

\begin{proof}[Proof of \cref{thm:critical-operator}]
Suppose that the critical connection kernel $T_{\pc}$ is unbounded.
Put $p_0=\pc/2$. For $p_0\le p<\pc$, use the corridor notation
\[
 \delta=\pc-p,\qquad S=\log(1/\delta),\qquad
 m_0=\norm{T_p},\qquad R=\frac{\chi_0^+}{m_0}\ge1,
 \qquad a_3=\tr(T_p^3).
\]
Sprinkling \eqref{eq:s05} with $a=\delta/(1-p)$ would bound
$T_{\pc}$ if $ad m_0<1$. The assumed unboundedness, together with
\eqref{eq:s22} and \eqref{eq:s23}, therefore gives
\begin{equation}
 m_0\ge\frac{1-p}{d\delta},\qquad
 \log R=o(S),\qquad \log a_3=o(S)
 \quad(\delta\downarrow0).
 \label{eq:s35}
\end{equation}
Indeed $m_0\ge c\delta^{-1}$, whereas, for every $\eta>0$,
$\chi_0^+\le C_\eta\delta^{-1-\eta}$ and
$1\le a_3\le C_\eta\delta^{-\eta}$. Hence
$1\le R\le C'_\eta\delta^{-\eta}$ for every $\eta>0$.

\begin{claim}[Comparison of deformed row sums]
\label{con:row-comparison}
There is $c_f>0$, depending only on $G$, such that
\begin{equation}
 s_t^V(z)\ge c_f\chi_t
 \qquad\text{for every $z\in V$ whenever $tc_k\le1$}.
 \label{eq:s36}
\end{equation}
\end{claim}

\begin{proof}
Connectedness and the finite number of vertex types give an integer
$L_f$ such that every vertex is within distance $L_f$ of every vertex
orbit. For fixed $t$, choose an orbit on which $s_t^V$ is maximal and
a simple path of length at most $L_f$ from $z$ to that orbit.
Equation \eqref{eq:s32}, with $p\ge p_0$, gives
\[
 s_t^V(z)\ge e^{-L_f}p_0^{L_f}\max_x s_t^V(x)
 \ge e^{-L_f}p_0^{L_f}\chi_t.
\]
Thus $c_f=e^{-L_f}p_0^{L_f}$ is admissible.
\end{proof}

Fix $0<\tau<c_f/(4c_\sigma)$ and the constants of
\cref{path:selection} for this value of $\tau$. They are now fixed
throughout the limit $p\uparrow\pc$. Since
$j\le J+1+C_1\log((k+1)R)$, there are fixed constants $c>0$ and
$C_2>0$ such that
\begin{equation}
 p^{2j}\ge c(kR)^{-C_2}\qquad(k\ge k_0,\ R\ge1).
 \label{eq:s37}
\end{equation}
For example, one may take
$C_2=2C_1\log(1/p_0)$ and
$c=p_0^{2(J+1)}2^{-C_2}$.

Put $\lambda=-\log\rho>0$ and choose a fixed $K_*$ with
$\kappa:=\lambda K_*>C_2+4$. Set
\[
 A_\delta=\log(2R)+\log(2+a_3)+\log S,\qquad
 k=\lceil K_*A_\delta\rceil,\qquad
 B_*=(kR)^{C_2+2},\qquad t_*=\frac{B_*}{m_0}.
\]
By \eqref{eq:s35}, $k\to\infty$ and $k=o(S)$; in particular
$k\ge k_0$ and $k\le S$ eventually. The estimates needed below are
\begin{equation}
 B_*\rho^k=o(1),\qquad
 R^2a_3\rho^k=o((kR)^{-C_2}),\qquad
 t_*jc_k=o(1).
 \label{eq:s38}
\end{equation}
To check their simultaneous validity, the definition of $k$ gives
\[
 \rho^k\le (2R)^{-\kappa}(2+a_3)^{-\kappa}S^{-\kappa}.
\]
Since $R,a_3\ge1$ and $k\le S$, it follows that
\[
 B_*\rho^k\le C S^{C_2+2-\kappa}\longrightarrow0,
 \qquad
 (kR)^{C_2}R^2a_3\rho^k
 \le C S^{C_2-\kappa}\longrightarrow0.
\]
Finally $c_k=e^{Ck}$ and $m_0\ge c e^S$, so
\[
 \log(t_*jc_k)
 \le -S+(C_2+2)(\log k+\log R)+\log j+Ck+O(1)
 =-S+o(S).
\]
This proves \eqref{eq:s38} with every path constant already fixed.

Consider the deformation with these values of $p,k$, on
$0\le t\le t_*$. Equation \eqref{eq:s27} and the first estimate in
\eqref{eq:s38} give
\begin{equation}
 m_t\ge m_0(1-2B_*\rho^k)\ge\frac{m_0}{2},\qquad
 \chi_t\ge\frac{m_0}{2c_\sigma}.
 \label{con:retained-norm}
\end{equation}
The second inequality uses the row-sum bound $m_t\le c_\sigma\chi_t$;
it requires no positivity of $T_t$ as an operator.

For a selected raw path in \cref{path:selection}, let $z$ be either
of its two escape endpoints. By \eqref{eq:s28}, \eqref{eq:s33},
\eqref{eq:s36}, and \eqref{con:retained-norm},
\[
 \begin{split}
 s_t^{D_w}(z)
 &\ge c_f\chi_t-\chi_0^+\frac{\tau}{R}
  =c_f\chi_t-\tau m_0\\
 &\ge(c_f-2c_\sigma\tau)\chi_t
  \ge\frac{c_f}{2}\chi_t.
 \end{split}
\]
The hypotheses of \eqref{eq:s36} hold throughout this interval by the
last estimate in \eqref{eq:s38}. Apply \eqref{eq:s32} to the two simple
escape paths in $D_w$, each of length at most $j$. Their product is
bounded below by
\[
 s_t^{D_w}(u)s_t^{D_w}(v)
 \ge e^{-2tjc_k}p^{2j}\left(\frac{c_f}{2}\chi_t\right)^2.
\]
Such paths have $\E_{\mathrm{path}}$-probability at least $b_*$, so
\eqref{eq:s37} and \eqref{eq:s38} imply
\begin{equation}
 \E_{\mathrm{path}}\!\left[
  \one_{\{\ell_w\ge3\}}s_t^{D_w}(u)s_t^{D_w}(v)\right]
 \ge c_1(kR)^{-C_2}\chi_t^2,
 \label{con:positive-term}
\end{equation}
uniformly on $[0,t_*]$, with $c_1>0$ independent of sufficiently small
$\delta$. On the other hand, the error in \eqref{eq:s31} is at most
\[
 C'(\chi_0^+)^2\rho^ka_3
 =C'R^2m_0^2\rho^ka_3
 \le4C'c_\sigma^2R^2\rho^ka_3\chi_t^2
 =o((kR)^{-C_2}\chi_t^2).
\]
The prefactor $e^{-tc_k}$ in that equation is uniformly bounded below.
We conclude that, for a constant $c_2>0$,
\begin{equation}
 -\chi_t'\ge c_2(kR)^{-C_2}\chi_t^2
 \qquad(0\le t\le t_*).
 \label{con:differential-decrease}
\end{equation}
The differentiability and domination established for the deformation
justify integration. Since $\chi_t>0$, it yields
\[
 \frac1{\chi_{t_*}}
 \ge c_2(kR)^{-C_2}t_*
 =\frac{c_2(kR)^2}{m_0}.
\]
But \eqref{con:retained-norm} gives
$1/\chi_{t_*}\le2c_\sigma/m_0$. These inequalities are incompatible
with $kR\to\infty$. Thus $T_{\pc}$ is bounded, proving
\cref{thm:critical-operator}.
\end{proof}

The next elementary endpoint lemma records the equivalence between
critical boundedness and the uniform subcritical formulation in
\cref{thm:main}.

\begin{lemma}[Monotone passage to an endpoint]\label{con:critical-limit}
Let $G$ be an infinite locally finite graph, let $p_*\in(0,1]$, and let
$T_p(x,y)=\Prob_p(x\leftrightarrow y)$ be its Bernoulli bond connection
kernels. Then $T_{p_*}$ is bounded on counting-measure $\ell^2(V)$ if
and only if $\sup_{0\le p<p_*}\norm{T_p}_{2\to2}<\infty$. In that case
the supremum equals $\norm{T_{p_*}}_{2\to2}$.
\end{lemma}

\begin{proof}
Each finite-domain connection probability is continuous in $p$, and
connectivity in $G$ is the increasing union of finite-domain
connectivity events. Thus $T_p(x,y)\uparrow T_{p_*}(x,y)$ as
$p\uparrow p_*$. If the subcritical norms are bounded by $M$, then
for finitely supported $f,g$ the finite sum
$\ip f{T_pg}$ tends to $\ip f{T_{p_*}g}$ and has absolute value at
most $M\norm f_2\norm g_2$. The limiting sesquilinear form extends
by density to a bounded operator with the required matrix entries
and norm at most $M$. Conversely, entrywise nonnegativity gives
\[
 |\ip f{T_pg}|\le\ip{|f|}{T_p|g|}
 \le\ip{|f|}{T_{p_*}|g|}
 \le\norm{T_{p_*}}\norm f_2\norm g_2.
\]
The two inequalities prove both equivalence and equality of norms.
\end{proof}

The simultaneous phase theorem of H\"aggstr\"om, Peres, and Schonmann
\cite[Theorem~1.3]{HaggstromPeresSchonmann1999} applies to all
quasi-transitive graphs and gives one event valid throughout the open
nonuniqueness phase. The following direct argument proves the compact
interval consequence from operator boundedness alone. It avoids taking
an uncountable intersection of events of probability one.

\begin{proposition}[One coupling for an interval]
\label{con:coupled-interval}
Let $G$ be infinite, connected, locally finite, and quasi-transitive.
Suppose $\pc<p_1<p_2<1$ and $T_{p_2}$ is bounded on $\ell^2(V)$.
In the coupling obtained by assigning independent uniform labels to
edges and opening an edge at parameter $p$ when its label is at most
$p$, almost surely there are infinitely many infinite clusters
simultaneously at every $p\in[p_1,p_2]$.
\end{proposition}

\begin{proof}
Let $N_{12}$ be the number of $p_2$-clusters containing an infinite
$p_1$-cluster. First, every automorphism-invariant event of the label
field has probability zero or one. Indeed, approximate such an event
$E$, within error $\varepsilon$ in probability, by a finite-coordinate
event $A$. Its coordinates lie in a finite ball centered at a vertex
in an unbounded orbit; such an orbit exists because there are finitely
many orbits and infinitely many vertices. An automorphism can move
this ball disjointly from itself. The events $A$ and its translate are
then independent, whereas $E$ agrees with its translate. It follows
that $\abs{\Prob(E)-\Prob(E)^2}\le4\varepsilon$, and letting
$\varepsilon\downarrow0$ proves the assertion.
In particular the invariant, measurable random variable
$N_{12}\in\{0,1,2,\ldots,\infty\}$ is constant almost surely.
Since $p_1>\pc$, an infinite $p_1$-cluster exists with positive
probability, so this constant is positive.

It cannot equal one. Write
$\theta(x)=\Prob_{p_1}(|C_x|=\infty)$. Some $\theta(x)$ is positive;
positive association with a finite path makes $\theta$ positive at
every vertex. Invariance and the finite number of types give
$\theta_*:=\inf_x\theta(x)>0$. If $N_{12}=1$, any two vertices in
infinite $p_1$-clusters belong to the same $p_2$-cluster. Positive
association therefore implies
\[
 T_{p_2}(x,y)\ge\theta(x)\theta(y)\ge\theta_*^2
 \qquad(x,y\in V).
\]
For finite sets $F$ of arbitrarily large size this gives
$\ip{\one_F}{T_{p_2}\one_F}/\norm{\one_F}_2^2
\ge\theta_*^2|F|$, contradicting boundedness.

Suppose instead that $N_{12}=n$ almost surely for an integer $n\ge2$.
By countability there are fixed vertices $x,y$ which, with positive
probability, belong to infinite $p_1$-clusters lying in distinct
$p_2$-clusters. Fix a finite path from $x$ to $y$. Replace its edge
labels by independent uniform labels in $[0,p_1]$. The resulting law
is absolutely continuous with respect to the original product law:
only finitely many coordinates were resampled, each with a bounded
density. At both parameters this modification only opens edges, and
on the specified positive-probability event it merges two counted
$p_2$-clusters.

Every infinite $p_1$-cluster after the modification contains an
infinite $p_1$-cluster from before it. Indeed finitely many added edges
can combine only finitely many old clusters into any new cluster; a
finite union of finite clusters is finite. Thus each new counted
$p_2$-cluster contains an old counted $p_2$-cluster, and old
$p_2$-clusters only merge. On the event in question the new count is
at most $n-1$. This contradicts absolute continuity and the almost
sure equality $N_{12}=n$. Hence $N_{12}=\infty$ almost surely.

On this single event, choose an infinite $p_1$-cluster in each of
infinitely many distinct counted $p_2$-clusters. For every
$p\in[p_1,p_2]$, their containing $p$-clusters are infinite and remain
distinct because their $p_2$-clusters are distinct. This proves the
simultaneous assertion.
\end{proof}

\begin{proof}[Proof of \cref{thm:main}]
First suppose $G$ is simple. If $\Aut(G)$ is unimodular,
\cref{thm:critical-operator} gives $\norm{T_{\pc}}<\infty$.
Since $\pc<1$, sprinkling \eqref{eq:s05} at $p=\pc$ with any
sufficiently small $a>0$ gives a bounded connection kernel at
$\pc+(1-\pc)a>\pc$. Thus $\pc<\ptwo$.

If $\Aut(G)$ is nonunimodular, \cref{pre:unimodular-reduction} supplies both
critical boundedness and the strict operator gap. In either case,
\cref{pre:threshold-order} yields $\pc<\ptwo\le\pu$.

If loops or parallel edges are allowed, use the simple subdivision
$H$ from \cref{pre:simple-reduction}. Infinite-cluster counts in $H$
at parameter $t$ agree with those in $G$ at parameter $t^2$, so
$\pc(H)^2=\pc(G)$. Moreover, on the old vertex set,
\[
 T^H_t\big|_{V(G)\times V(G)}=T^G_{t^2}.
\]
Compression of a bounded operator to this coordinate subspace is
bounded. Applying the simple-graph result to $H$ therefore gives
boundedness of $T^G_{\pc(G)}$ and of $T^G_p$ at some $p>\pc(G)$.
The inequality $\ptwo(G)\le\pu(G)$ again follows from
\cref{pre:threshold-order}. This covers the full graph class in the
statement.

Finally choose $\pc<p_1<p_2<1$ with $T_{p_2}$ bounded, which the strict
operator gap permits. \Cref{con:coupled-interval} supplies
the claimed simultaneous nonempty interval of infinitely many
infinite clusters. The equality of the uniform subcritical and critical
norms is \cref{con:critical-limit}, completing the proof.
\end{proof}
